\documentclass[conference]{IEEEtran}
\IEEEoverridecommandlockouts

\usepackage{cite}
\usepackage{longtable}
\usepackage{booktabs}
\usepackage{array}
\usepackage{float}
\usepackage{graphicx}
\usepackage{textcomp}
\usepackage{xcolor}
\usepackage{amsmath,amssymb,amsfonts}
\usepackage{algorithm}
\usepackage{algpseudocode}
\usepackage{hyperref}

\def\BibTeX{{\rm B\kern-.05em{\sc i\kern-.025em b}\kern-.08em
    T\kern-.1667em\lower.7ex\hbox{E}\kern-.125emX}}

\newcommand{\tightlist}{\setlength{\itemsep}{0pt}\setlength{\parskip}{0pt}}

\begin{document}

\title{Single-Machine Optimization of Perturbation-Based Loss-Grid Computation\\
\thanks{This thesis was prepared in partial fulfilment of the requirements for the Degree of Bachelor of Computer Science, Maastricht University. Supervisor(s): Max Sondag, Georgios Stamoulis}
}

\author{\IEEEauthorblockN{Given Name Surname}
\IEEEauthorblockA{\textit{Department of Advanced Computing Sciences} \\
\textit{Faculty of Science and Engineering}\\
\textit{Maastricht University}\\
Maastricht, The Netherlands}
}

\maketitle

\begin{abstract}
TODO left for last
\end{abstract}


\section{INTRODUCTION}\label{introduction}

\subsection{Motivation and Scope}\label{motivation-and-scope}

A loss landscape is a 2D
scalar field made by perturbing a trained parameter vector
along two directions and evaluating the loss at each
point~\cite{li2018losslandscape}, supporting
generalization analysis and interactive
checkpoint-comparison tools~\cite{xie2025losslens,losslens_repo}.

Computationally, the landscape is represented as an array of
scalar loss evaluations indexed by grid coordinates. Each entry
evaluates the model at the corresponding perturbed parameter
vector on data. One landscape can require hundreds to thousands
of such
evaluations~\cite{li2018losslandscape,goldstein_loss_landscape_repo}.
In interactive checkpoint comparison, this repeated work can
dominate latency on commodity hardware.

\subsection{Problem Statement}\label{problem-statement}

The problem is to evaluate candidate single-machine optimization
mechanisms for perturbation-based loss-grid computation under the
constraint of one GPU, the available CPU cores, and finite RAM
and VRAM. The optimization target is the construction of the
loss grid itself: the perturbed model evaluations that fill the
grid. The thesis studies this target through two mechanisms:
faster grid-point evaluation on one device, and CPU/GPU
distribution on one machine. Because both mechanisms depend on
hardware and workload, the study determines where each produces
statistically supported runtime reductions, where it regresses or
becomes impractical, and whether it preserves the numerical
equivalence of the resulting loss surface.

\subsection{Research Questions}\label{research-questions}

\texttt{RQ1.} Do PyTorch's stateless evaluation APIs reduce
perturbation-based loss-grid runtime, and for which architecture types?

\texttt{RQ2.} Under what conditions do idle CPU cores reduce total
loss-grid time, and how is the CPU/GPU inference throughput ratio
associated with those conditions?

\texttt{RQ3.} Does reusing the hybrid scheduler's calibrated
policy across a related checkpoint session amortize the
upfront calibration cost?

\subsection{Structure of the Report}\label{structure-of-the-report}

Section~\ref{related-work} frames the thesis with respect to prior work.
Section~\ref{methods} defines the baseline and the three
optimization approaches (one per RQ).
Section~\ref{experiments} gives the experimental design,
workloads, platforms, and measurement protocol.
Section~\ref{results} reports measurements;
Section~\ref{discussion} interprets them as applicability
conditions; Section~\ref{conclusions} answers the RQs.

\subsection{Contributions}\label{contributions}

(RQ1)~Whether the \texttt{functional\_call} + \texttt{vmap}
pattern from PyTorch's model-ensembling
tutorial~\cite{pytorch_ensembling} transfers to loss-grid
construction --- a workload with the same outer-loop shape but
different semantics (one model, many perturbations) --- across
architectures and backends.
(RQ2)~The conditions under which idle CPU cores reduce grid
time, tied to the CPU/GPU inference throughput ratio that
Gallet and Gowanlock~\cite{gallet2021heterogeneous} identified
as the core variable of heterogeneous CPU/GPU scheduling
applicability for a different workload (set operations) and
that this thesis tests in the ML-inference regime.
(RQ3)~Whether reusing the hybrid scheduler's calibrated policy
across $N=4$ related checkpoints, motivated by LossLens-style
checkpoint comparison, amortizes the upfront calibration cost
--- measured as net cumulative session speedup, with the
break-even session size reported as a diagnostic. Cases
favoring the baseline are reported as primary findings
alongside the optimization wins.


\section{RELATED WORK}\label{related-work}

Single-machine loss-grid optimization intersects loss-landscape
visualization, adaptive loss-grid sampling, heterogeneous
CPU/GPU scheduling, and autotune-and-cache patterns.

\subsection{Loss-landscape visualization and
analytics}\label{rw-losslandscape}

Goodfellow et al.~\cite{goodfellow2015qualitative} introduced
1D linear interpolation between parameter points to
qualitatively characterize neural-network optimization
trajectories. Li et al.~\cite{li2018losslandscape} extended this
to a 2D loss grid with filter-normalized perturbation directions,
the algorithm taken here as the baseline; the official
implementation~\cite{goldstein_loss_landscape_repo} runs it under
MPI on a multi-node cluster. LossLens~\cite{xie2025losslens,losslens_repo}
integrated this construct into an interactive analytics
workflow over model modes/checkpoints.

\subsection{Adaptive and progressive loss-grid
sampling}\label{rw-adaptive-sampling}

Another line of work replaces the canonical full grid with
an adaptive sampler that evaluates loss only at points predicted
to carry visual information. Prakash et al.
apply this to loss-landscape visualization directly, replacing
regular $\alpha,\beta$ sampling with an adaptive algorithm
focused on ``interesting'' regions and comparing it against a
GradVis-style baseline and a vectorized uniform
sampler~\cite{prakash2024adaptive}. Their reported outcome is
mixed: adaptive sampling beats the GradVis-style baseline but
loses to the vectorized uniform sampler due to selection
overhead, and their user-preference comparison is
inconclusive~\cite{prakash2024adaptive}.

\subsection{Heterogeneous compute and ML inference
throughput}\label{rw-scheduling}

Task-parallel heterogeneous scheduling assigns whole
independent tasks across devices. Dynamic self-scheduling,
introduced by Polychronopoulos and
Kuck~\cite{polychronopoulos1987gss}, is the simplest variant:
idle workers poll a shared queue and claim one task at a time,
requiring no per-device performance model. Richer policies in
the same family include HEFT-style cost-model scheduling and
priority-based work-stealing, exemplified by
StarPU~\cite{augonnet2011starpu}; these address load
imbalance and per-task cost variance. Gallet and
Gowanlock~\cite{gallet2021heterogeneous} identify the
CPU/GPU throughput ratio as an applicability variable for
heterogeneous task scheduling, demonstrated on a CPU--GPU
epsilon grid join: static and dynamic partitioning yield a
speedup over GPU-only execution only above a workload-dependent
ratio threshold, and regress below it.

An orthogonal approach splits a single computation across
devices, with per-device stages exchanging intermediate data
along a pipeline. OpenVINO's HETERO execution mode is a
production example at the intersection of heterogeneous
compute and ML inference~\cite{openvino_hetero}: with
pipeline-parallel \texttt{model\_distribution\_policy}, the
inference model is divided into stages, each stage is assigned
to a device, and the output of one stage is fed as input to
the next.

Production ML inference systems also pursue throughput
through request-level batching. OpenVINO's automatic
batching~\cite{openvino_auto_batching,openvino_auto_batching_src}
groups concurrent batch-1 inference requests at runtime to
improve device utilization on supported devices and
throughput-oriented modes. At implementation level,
auto-batching operates on the input/request batch dimension
of one fixed compiled model: requests are packed into slices
of a larger input/output tensor and executed by a reshaped
batched model.

\subsection{Autotune-and-cache patterns}\label{rw-autotuning}

ATLAS~\cite{whaley1998atlas} introduces an autotune-and-cache
pattern: empirically tune the inner loops of BLAS routines
(GEMM, DAXPY, etc.) for the target machine's cache and
register hierarchy; cache the selected configuration by
(machine, BLAS-op) signature; reuse it across every
subsequent matching call. The pattern is four-part: tune
once per machine, cache by signature, reuse across matching
invocations, per-machine scoping. OpenTuner~\cite{ansel2014opentuner}
canonized bounded autotuning search with explicit stopping
criteria as the discipline-level departure from ATLAS-style
exhaustive search.

\subsection{Where this thesis sits}\label{rw-placement}

Prior work provides the core
algorithm~\cite{li2018losslandscape,goldstein_loss_landscape_repo}
and its interactive checkpoint-comparison use
case~\cite{xie2025losslens,losslens_repo}, the dynamic
self-scheduling
primitive~\cite{polychronopoulos1987gss} and the broader
heterocompute-scheduler family~\cite{augonnet2011starpu},
the throughput-ratio applicability
variable~\cite{gallet2021heterogeneous} from a different
workload, the production model-parallel and request-batching
counterparts in ML
inference~\cite{openvino_hetero,openvino_auto_batching}, the
autotune-and-cache pattern~\cite{whaley1998atlas} with its
bounded-search discipline~\cite{ansel2014opentuner}, and the
adjacent adaptive-sampling
direction~\cite{prakash2024adaptive}.

The gap is the combination of workload, deployment target, and
optimization question. Existing loss-landscape work defines the
canonical grid and motivates interactive checkpoint comparison,
but does not evaluate commodity single-machine optimization of
the grid-construction runtime. Adaptive samplers change which
grid points are evaluated; this thesis keeps the canonical grid
fixed. Heterogeneous task-scheduling work supplies the
CPU/GPU-throughput-ratio variable, but from a different workload
with different task granularity and kernels. Production ML
inference systems are closer in domain, but optimize adjacent
axes: HETERO partitions one compiled model's operation graph
across devices, automatic batching groups concurrent inference
requests along the input/request batch dimension, and PyTorch's
\texttt{functional\_call}+\texttt{vmap} pattern vectorizes a
stack of parameter states on one device. None of these directly
targets the loss-grid workload shape: one model family, many
perturbations of one checkpoint, shared data, and finite CPU,
GPU, RAM, and VRAM on a commodity machine.



\section{METHODS}\label{methods}

This section defines the canonical loss-grid algorithm and the
three optimization approaches mapped to the research questions:
algorithm-level optimization candidates
(Section~\ref{methods-redesign}, RQ1), multi-device
heterogeneous scheduling (Section~\ref{methods-scheduling},
RQ2), and adaptive calibration with cache reuse
(Section~\ref{methods-calibration}, RQ3). It also states the
shared benchmarking, surface-validation, and uncertainty
protocol. Terminology is defined in
Section~\ref{introduction}; experimental procedures in
Section~\ref{experiments}.

\subsection{Canonical Loss-Grid Algorithm
(Baseline)}\label{methods-canonical-algorithm}

The thesis takes Li et al., \emph{Visualizing the Loss
Landscape of Neural Nets}~\cite{li2018losslandscape}, as the
canonical method. It is the de-facto reference in subsequent
work~\cite{xie2025losslens,losslens_repo} and the algorithm
shipped to end-users~\cite{goldstein_loss_landscape_repo};
alternative formulations (e.g.\ trajectory sampling) probe a
different aspect and are not the cost driver in
LossLens-style workflows.

For one grid point the algorithm picks coordinates
$(\alpha,\beta)$, perturbs parameters via $w(\alpha,\beta) =
w_0 + \alpha\,d_1 + \beta\,d_2$, evaluates loss over the
dataset subset, and writes the scalar into the surface; this
repeats until the grid is complete. Here $w_0$ is the trained
parameter vector and $d_1, d_2$ are filter-normalized
directions following the upstream default \texttt{weights}
mode~\cite{li2018losslandscape,goldstein_loss_landscape_repo}.

The thesis adopts the procedure as a controlled single-machine
comparison baseline with two adjustments to align with the
research questions. First, the baseline uses one GPU, matching
the problem framing in Section~\ref{problem-statement} and
leaving RQ2 to measure CPU--GPU scheduling rather than
GPU--GPU distribution. Second, the baseline performs no CPU
helper execution, because hybrid execution is itself the RQ2
candidate; including it in the baseline would make that
comparison circular. Because the baseline is single-GPU, the
implementation runs in a single process rather than through
the public MPI/distributed path. The only omission relative
to the public MPI
reference~\cite{goldstein_loss_landscape_repo} is distributed
execution. The per-grid-point semantics are preserved, but
grid points are evaluated sequentially. All comparisons in
Section~\ref{results} are relative to this baseline. A
perturbed setting $w(\alpha,\beta)$ is a \emph{perturbation}
of a checkpoint; \emph{checkpoint} is reserved for a distinct
trained model from a separate run.

\subsection{Algorithm-Level Optimization Candidates
(RQ1)}\label{methods-redesign}

The canonical loop evaluates one perturbed parameter state per
grid point against the same data subset. This is the same
outer-loop shape that PyTorch's model-ensembling
tutorial~\cite{pytorch_ensembling} addresses for stacks of
independently trained models, via the composition of
\texttt{functional\_call}~\cite{pytorch_functional_call} and
\texttt{vmap}~\cite{pytorch_vmap}.
\texttt{torch.compile}~\cite{pytorch_compile} is an orthogonal
intra-device mechanism — ahead-of-time graph capture, kernel
fusion, and dispatch elision — on the same per-grid-point
work. RQ1 tests three candidates against the baseline
(Table~\ref{tab:opt-map}): \texttt{vmap} composed with
\texttt{functional\_call} (the documented ensembling pattern),
\texttt{torch.compile} on the baseline evaluation path, and
the composition of both.
\texttt{functional\_call} is a prerequisite primitive rather
than a standalone candidate: it makes no independent
optimization claim, and the candidates that build on it
subsume its observable runtime contribution.

The per-candidate verdict is the effect-size CI on the
$T_{\text{grid}}$ speedup ratio
(Section~\ref{methods-benchmarking}). It is paired with an
\emph{intra-axis composition verdict}
$\in \{\text{multiplicative}, \text{redundant},
\text{destructive}, \text{inconclusive}\}$, derived by
comparing \texttt{compiled\_vmapped}'s observed speedup to the
product of \texttt{vmapped}'s and \texttt{compiled}'s
individual speedups under the same effect-size CI rule on the
ratio.

\begin{table}[t]
\caption{Algorithm-level optimization candidates.}
\label{tab:opt-map}
\centering
\small
\begin{tabular}{@{}p{0.27\linewidth}p{0.27\linewidth}p{0.36\linewidth}@{}}
\toprule
Candidate & Components & Method role \\
\midrule
baseline & in-place mutation + sequential loop & reference implementation \\[4pt]
\texttt{vmapped} & \texttt{functional\_call} + \texttt{vmap} + chunking & vectorized forward/loss over perturbations \\[4pt]
\texttt{compiled} & \texttt{torch.compile} on baseline eval path & graph capture and kernel fusion on per-grid-point work \\[4pt]
\texttt{compiled\_vmapped} & \texttt{torch.compile} over the \texttt{vmapped} candidate & intra-axis composition: graph capture over vectorized point evaluation \\
\bottomrule
\end{tabular}
\end{table}

\subsubsection{\texorpdfstring{Vectorized Point Evaluation with
\texttt{torch.vmap}}{Vectorized Point Evaluation with torch.vmap}}\label{vectorized-point-evaluation-with-torch.vmap}

Loss-grid computation has the same outer-loop shape as model
ensembling: a stack of parameter states evaluated against the
same data. The model-ensembling
tutorial~\cite{pytorch_ensembling} composes
\texttt{vmap}~\cite{pytorch_vmap} with
\texttt{functional\_call}~\cite{pytorch_functional_call} to
evaluate that stack in one batched call.
\texttt{functional\_call} replaces in-place parameter mutation
with a stateless invocation: the perturbed parameter mapping
is supplied to the model call for that invocation, leaving
the module object unrewritten. \texttt{vmap} then applies the
loss computation across a leading batch dimension over the
stacked parameter tensors. The \texttt{vmapped} candidate
instantiates this composition: it stacks the perturbed
parameter tensors for $K$ grid points into batched tensors of
shape $(K,S)$ and applies the loss computation across the $K$
slices per data batch
(Figure~\ref{fig:point-evaluation-paths}).

\begin{figure}[t]
\centering
\begin{minipage}{0.96\linewidth}
\small
\textbf{Functional parameter application for grid point $i$}
\begin{algorithmic}[1]
\Statex \textbf{Baseline path}
\State $\mathit{model}.\mathit{parameters} \gets \mathit{params}_i$
\State $\mathit{loss}_i \gets \mathit{model}(\mathit{batch})$
\Statex
\Statex \textbf{Functional path}
\State $f_i \gets \textsc{FunctionalCall}(\mathit{model}, \mathit{params}_i)$
\State $\mathit{loss}_i \gets f_i(\mathit{batch})$
\end{algorithmic}

\vspace{0.6em}

\textbf{Vectorized point evaluation over the grid}
\begin{algorithmic}[1]
\Statex \textbf{Baseline path}
\For{$\mathit{point} \in \mathit{grid}$}
  \For{$\mathit{batch} \in \mathit{data}$}
    \State $\mathit{loss}(\mathit{point}, \mathit{batch})$
  \EndFor
\EndFor
\Statex
\Statex \textbf{Vmapped path}
\For{$\mathit{point\_chunk} \in \mathit{grid\_chunks}$}
  \For{$\mathit{batch} \in \mathit{data}$}
    \State $\mathit{losses}(\mathit{point\_chunk}, \mathit{batch})$
  \EndFor
\EndFor
\end{algorithmic}
\end{minipage}
\caption{Baseline, functional, and vmapped point-evaluation paths.}
\label{fig:point-evaluation-paths}
\end{figure}

The predicted mechanism, if the pattern transfers, is dispatch
amortization and improved use of backend batching rules;
actual kernel behavior is backend- and architecture-dependent.
Full-grid vectorization may exceed memory, so the candidate
evaluates chunks of $K$ grid points at a time and sweeps $K$
for breadth.

\subsubsection{\texorpdfstring{Graph Capture with
\texttt{torch.compile}}{Graph Capture with torch.compile}}\label{graph-capture-with-torch.compile}

\texttt{torch.compile}~\cite{pytorch_compile} captures the
forward/loss path as an ahead-of-time graph, fuses kernels
where the backend permits, and elides Python-level dispatch.
The \texttt{compiled} candidate applies \texttt{torch.compile}
to the baseline per-grid-point evaluation path, isolating the
graph-capture/fusion mechanism from any batching effect; the
loop structure is unchanged. The \texttt{compiled\_vmapped}
candidate applies \texttt{torch.compile} to the
\texttt{vmapped} candidate's batched function, testing
whether graph capture and \texttt{vmap}'s batching multiply on
the same work or interfere. PyTorch documents no canonical
composition of \texttt{torch.compile} with \texttt{vmap} on
the loss-grid pattern, so the composition verdict is
empirical.

\subsection{Multi-Device Heterogeneous Scheduling
(RQ2)}\label{methods-scheduling}

The second approach evaluates heterogeneous CPU/GPU scheduling
for parallel grid-point evaluation. Its premise inherits the
CPU/GPU throughput ratio as the core variable of
heterogeneous CPU/GPU scheduling applicability, identified by
Gallet and Gowanlock~\cite{gallet2021heterogeneous} for set
operations: near or above parity, the CPU can absorb a
non-trivial share of tasks; far below parity, the GPU-only
baseline is already near-optimal. The scheduling mechanism is
dynamic self-scheduling of Polychronopoulos and
Kuck~\cite{polychronopoulos1987gss}, the simpler variant of
the StarPU class of heterocompute
schedulers~\cite{augonnet2011starpu}. The simpler variant is
sufficient because loss-grid is embarrassingly parallel with
uniform per-task work, removing the load-imbalance and
heterogeneous-task-cost problems that motivate the richer
StarPU policies.

The scheduler is eager: GPU and CPU workers poll a shared
queue, one grid point per task, with the same
surface-validation contract as the baseline. Both devices run
the same forward+loss code, keeping the CPU/GPU throughput
ratio interpretable as RQ2's predictor. CPU workers are
bounded by the available logical cores. One grid point per
task is the natural unit of the embarrassingly-parallel
algorithm; no aggregation is needed because per-point work is
uniform.

Real hardware fixes the native ratio, so the method exposes
an optional \texttt{gpu\_slowdown\_factor} that injects a
controlled post-kernel delay on the GPU worker, moving the
operating point along the ratio axis so the
ratio-versus-speedup relationship can be measured across
regimes the native workload set does not span. This is the
software analog of the GPU-throughput-sweep methodology
established by Mei and Chu~\cite{mei2017gpudvfs} (where the
native knob is DVFS); the sweep instrument is methodological,
not a scheduling mechanism. Native and achieved CPU/GPU
ratios are reported alongside each measurement.

The delay-injection knob is not a competing approach to
hardware-faithful GPU characterization. Validated
simulators~\cite{khairy2020accelsim}, kernel
sampling~\cite{avalos2021pka}, and model-driven throughput
predictions~\cite{gabbay2022dnnmemory} attribute runtime to
architectural causes inside the device; here the GPU is
treated as a black box and a controllable offset is added
to traverse the ratio axis without modeling kernel-level
contention. The purpose is to reach regimes the native
workload set does not span, not to explain why a given
throughput is observed.

\subsection{Adaptive Calibration with Cache Reuse
(RQ3)}\label{methods-calibration}

The hybrid scheduler (Section~\ref{methods-scheduling})
requires a per-machine calibration step to select its
hetero-config tuple: a policy choice $\in
\{\texttt{gpu\_only}, \texttt{gpu\_cpu\_hybrid}\}$ together
with GPU batch size, CPU worker count, and CPU batch size
when hybrid is selected. This calibration is the
unconditional upfront cost paid before any per-checkpoint
savings can be realized. A single-checkpoint user pays that
cost without reusing the selected tuple and loses; the cache
changes the unit of cost from per-checkpoint to per-session by
paying calibration once and reusing the selected tuple across
$N$ related checkpoints on the same machine. Loss-landscape
analysis is
rarely a single run: users compare related
checkpoints~\cite{xie2025losslens,losslens_repo}. A
\emph{related checkpoint} shares architecture, dataset, loss,
and grid spec but differs in learned weights (e.g.\ a
different seed).

\textbf{Composition with RQ1.} The GPU-side implementation
inside the hetero-config tuple is configured as RQ1's
per-(machine, workload) winner for the same workload; RQ1's
choices apply to the GPU worker only, since \texttt{vmap}-
and \texttt{torch.compile}-style transforms over the small
per-task CPU slices are out of scope. RQ1 is not in RQ3's
sweep — its winner is established by RQ1 and consumed here as
a configuration input. Composition with RQ1 happens at the
deployed-stack level, not the sweep level: the session
runtime assembles RQ1's winner with RQ3's calibrated
hetero-config tuple. The cache key (machine, workload)
matches RQ1's key so the composition is well-defined and the
amortization claim is keyed on the same identity.

\textbf{Calibration probe.} Calibration uses the smallest
square grid satisfying $r^2 \geq 4 \times (1 + p_{\max})$,
where $1 + p_{\max}$ is the total compute-unit count (one
GPU plus $p_{\max}$ CPU workers). The factor $4$ keeps the
per-unit task count large enough that scheduler granularity
is not the limiting factor of the probe. The probe sweeps
GPU batch size, CPU worker count, and CPU batch size with
patience-based
termination after \texttt{calibration\_retry}$=R$
consecutive non-improvements, inheriting OpenTuner's bounded
autotuning-search discipline~\cite{ansel2014opentuner} as the
departure from ATLAS-style exhaustive search. The selected
configuration is the lowest-runtime valid one, or
\texttt{gpu\_only} fallback if no hybrid configuration beats
the GPU-only baseline.

\textbf{Cache reuse.} For a related-checkpoint session, the
selected tuple is plausibly stable across checkpoints, so
the calibration cost paid by the first checkpoint is reused
for the rest. This instantiates the autotune-and-cache
pattern of ATLAS~\cite{whaley1998atlas} with the signature
lifted from op-level to workload-level: ATLAS keys on
(machine, BLAS op), RQ3 keys on (machine, workload =
architecture + dataset + loss + grid spec). The cached value
is the hetero-config tuple selection. Cache reuse is
invalid if the selected mode is not stable across the
session.

\textbf{Session comparison.} RQ3 is measured in the native
operating regime of the machine. The calibrated tuple is used
for every checkpoint in the composed session, and the
reference session runs the vanilla per-checkpoint path. The
reported quantities are the one-time calibration cost, the
cached per-checkpoint grid time, cumulative
$T_{\text{session}}$, and the break-even session size.

\subsection{Benchmarking, Surface Validation, and
Uncertainty}\label{methods-benchmarking}

\textbf{Timing primitive.} Elapsed time is measured at the
caller with \texttt{time.perf\_counter\_ns} as specified by
PEP 418~\cite{python_perf_counter_ns}: a portable monotonic
high-resolution clock that captures host-observed wall time
including Python orchestration, framework dispatch, backend
execution, parameter binding, dataloader iteration, result
collection, and synchronization. Asynchronous backend
completion is forced through the framework backend
abstraction before stopping the timer; no
CUDA-event-specific dependency is introduced. PyTorch Timer
and CUDA events are not root canon for this thesis.

\textbf{Runtime targets.} Two targets are named separately:
$T_{\text{grid}}$ is the wall time for one loss-grid after
inputs are ready (the RQ1 and RQ2 target), and
$T_{\text{session}}$ is the wall time for a multi-checkpoint
session, including any one-time calibration the measured arm
performs (the RQ3 target). Every speedup claim names the
target it uses. $T_{\text{grid}}$ is also reported in a
section decomposition $T_{\text{grid}} = T_{\text{perturb}} +
T_{\text{bind}} + T_{\text{eval}} + T_{\text{collect}} +
\text{residual}$ using the same wall-time primitive on each
section. The decomposition is explanatory only; it is not a
separate timing mechanism.

\textbf{Repeats and uncertainty.} Repetition and stopping
budget follow Kalibera and
Jones~\cite{kalibera2013rigorous}: repeated runs after
warm-up, variation estimates per arm, and a stopping budget
that bounds total experiment time. The speedup ratio is
reported with an effect-size confidence interval following
Kalibera and Jones~\cite{kalibera2020effect}. The verdict
rule is: $\text{CI}_{\text{low}} > 1.0 \Rightarrow$
supported speedup; $\text{CI}_{\text{high}} < 1.0 \Rightarrow$
supported regression; otherwise inconclusive. The same rule
applies across RQ1, RQ2, and — when full sessions are
repeated — RQ3; RQ3 is descriptive session evidence
otherwise.

\textbf{Trial order and pairing.} Baseline and candidate are
evaluated within the same repeat under identical workload,
seed, and grid; trial order rotates across repeats to avoid
fixed-order benchmarks. The pairing scheme is disclosed per
experiment.

\textbf{Surface-validation gate.} A candidate's timing claim
is conditional on producing the same loss surface as the
baseline on the same (workload $\times$ backend $\times$
checkpoint). Different implementations of the same nominal
float32 computation produce different numeric outputs in
practice (op reordering, kernel/fusion choice, hybrid
CPU/GPU split), so the gate is a relative-equivalence
criterion rather than bitwise equality. Three parts:
\begin{enumerate}
\item grid-point count and coordinates match the baseline
exactly;
\item paired NaNs at matching coordinates count as agreement;
\item finite losses satisfy $|a-b| \leq 1\mathrm{e}{-5}
\cdot \max(|a|,|b|)$ — $\text{rtol} = 1\mathrm{e}{-5}$
follows the \texttt{torch.allclose}
default~\cite{pytorch_allclose}; $\text{atol} = 0$, so
equivalence requires proportionally small drift rather than
an absolute floor.
\end{enumerate}
This gate is stricter than the visual-equivalence convention
used by canonical loss-landscape
implementations~\cite{li2018losslandscape,goldstein_loss_landscape_repo},
which validate landscape topology by inspection rather than
by per-point tolerance.
A surface that fails the gate disqualifies the timing claim
regardless of measured speedup.


\section{EXPERIMENTS}\label{experiments}

This section states the workload set, platforms, and
experiment-specific design choices. Shared timing, uncertainty,
and surface-validation rules are defined once in
Section~\ref{methods-benchmarking}. The three experiments form
a decision pipeline. RQ1 decides which GPU evaluation mode to
use and exports it as \texttt{rq3\_config}. RQ2 decides
whether heterogeneous scheduling is worth considering and on
which workload/platform regimes yield positive evidence. RQ3
decides whether calibration reuse amortizes only for regimes
with positive evidence in RQ2. Each stage depends on the
previous result, so a null or negative result upstream narrows
the choices available downstream
(Sections~\ref{experiment-1-algorithm}--\ref{experiment-3-cache}).

\subsection{Shared Experimental
Setup}\label{workload-set-and-platforms}

\textbf{Platforms.} Table~\ref{tab:platforms} gives the two
hardware settings: Apple~M4 with unified memory and ten CPU
cores, and NVIDIA~Tesla~T4 with discrete VRAM, higher GPU
throughput, and two CPU cores. The unified-vs-discrete
memory axis and the CPU-core count separate the RQ2 operating
points and bound the vmap chunk memory budget in RQ1.

\begin{table}[t]
  \caption{Platform Characterization}
  \label{tab:platforms}
  \centering
  \begin{tabular}{lll}
    \toprule
    Property & Apple M4 & NVIDIA Tesla T4 \\
    \midrule
    GPU & M4 integrated & Tesla T4 \\
    VRAM & 16\,GiB (unified) & 14.6\,GiB (discrete) \\
    CPU cores available & 10 & 2 \\
    PyTorch backend & \texttt{mps} & \texttt{cuda} \\
    \bottomrule
  \end{tabular}
\end{table}

\textbf{Workload set.} A workload is a (dataset, model family, loss function) triple. Four workloads cover
distinct computation patterns relevant to the RQ1 transform
applicability claim (Table~\ref{tab:workloads}):
\texttt{cifar10\_resnet20} for 2D convolution (the
workload for which CUDA convolution kernels are
designed~\cite{nvidia_cudnn_convolution_guide}),
\texttt{cifar10\_row\_gru}~\cite{cho2014gru,chung2014gru}
for recurrent row dependencies (no kernel fusion),
\texttt{mnist\_mlp} for MLP image classification, and
\texttt{california\_mlp} for
tabular regression. The MLP pair differs in input dimensionality and
per-sample compute cost, the axes that should impact vmap behavior. 
These span enough architectures
for workload-dependent outcomes to emerge. RQ1 and RQ2 use
all four; RQ3 narrows to one case selected by filtering
Experiment~2 for hybrid affinity
(Section~\ref{experiment-3-cache}).

\begin{table}[t]
  \caption{Workload set.}
  \label{tab:workloads}
  \centering
  \begin{tabular}{llll}
    \toprule
    Name & Modality & Architecture & Dataset \\
    \midrule
    \texttt{cifar10\_resnet20} & image & 2D conv (ResNet-20) & CIFAR-10 \\
    \texttt{cifar10\_row\_gru} & sequence & row-wise GRU & CIFAR-10 \\
    \texttt{mnist\_mlp} & dense image & MLP & MNIST \\
    \texttt{california\_mlp} & tabular & MLP & California Housing \\
    \bottomrule
  \end{tabular}
\end{table}

\textbf{Grid and data subset.} RQ1 and RQ2 use an
$8\times 8$ filter-normalized grid at unit perturbation
scale, 1024 evaluation samples per point, and batch size
$64$. Unit scale matches the upstream default for
filter-normalized
weights~\cite{li2018losslandscape,goldstein_loss_landscape_repo}.
The $8\times 8$ resolution yields 64 grid-point tasks per
run: enough to populate the $K\in\{32,64\}$ vmap chunks in
RQ1 and to fill the shared queue in RQ2, while keeping the
full sweep feasible on both platforms. Resolution sets the
task count, not the per-point evaluation semantics, so the
transform and scheduler verdicts do not depend on it. RQ3 uses a
$20\times 20$ session grid matching LossLens's default
resolution~\cite{xie2025losslens}, plus a calibration probe
grid of side $m$ with $m^2\ge 4(1+p_{\max})$
(Section~\ref{methods-calibration}).

\textbf{Repeats and trial order.} RQ1 and RQ2 use $R=5$
paired repeats ($df=R-1=4$, the minimum for the paired
t-CI). The paired design removes machine variance between runs,
so a small $R$ resolves the per-candidate verdicts across the
full $4\times 2$ workload/platform sweep at feasible total
cost; larger $R$ tightens the CI with diminishing
returns~\cite{kalibera2013rigorous}. Within each repeat the candidates run sequentially
under a cyclic rotation indexed by the repeat number, so no
candidate consistently runs after baseline and inherits its
warm cache
state~\cite{kalibera2013rigorous}. RQ3 takes a single
$T_{\text{session}}$ measurement per (workload, platform,
condition) and is descriptive
(Section~\ref{methods-benchmarking}).

\textbf{Conditions.} \texttt{vanilla} denotes the canonical
GPU-only path defined in
Section~\ref{methods-canonical-algorithm}. RQ1 pairs each
candidate from Section~\ref{methods-redesign} against
\texttt{vanilla}; RQ2 pairs \texttt{hybrid} against
\texttt{vanilla} under matched GPU slowdown; RQ3 pairs
\texttt{cached\_composed} against \texttt{vanilla} run
once per checkpoint.

\subsection{Algorithm-Level Optimization
(RQ1)}\label{experiment-1-algorithm}

\textbf{Question.} Do PyTorch's stateless evaluation
primitives reduce $T_{\text{grid}}$, and which workload
patterns gain or regress?

\textbf{Hypothesis.} Verdicts depend on architecture. We
predict supported improvement on the dense MLP rows (the
closest match to the \texttt{vmap} ensembling pattern), an
unresolved or smaller effect on \texttt{cifar10\_resnet20}
(its convolution kernels already saturate the GPU at this
batch size), and an unresolved or supported regression on
\texttt{cifar10\_row\_gru} (its inner row-by-row dependence
is not removable by batching perturbations).

\textbf{Design.} The candidates from
Section~\ref{methods-redesign} run on each workload/platform
pair under the shared setup. \texttt{vmap}-bearing
candidates sweep $K\in\{32,64\}$ as a memory-vs-throughput
sweep: the smaller chunk bounds memory and isolates the
dispatch-amortization gain, the larger chunk tests whether
higher per-chunk throughput exceeds the added memory
pressure. The composed candidate is compared against the
better individual transform via $q_{\text{compose}}$
(Section~\ref{methods-redesign}).

\subsection{Hybrid Scheduler Applicability
(RQ2)}\label{experiment-2-hybrid}

\textbf{Question.} Does heterogeneous CPU+GPU scheduling
reduce $T_{\text{grid}}$, and does the CPU/GPU forward+loss
throughput ratio $r$ determine where it applies?

\textbf{Hypothesis.} Near parity, dynamic self-scheduling~\cite{kruskal1985selfscheduling}
absorbs useful CPU work and reduces $T_{\text{grid}}$
against the matched GPU-only path. Far below parity,
GPU-only dominates or the result is unresolved. We also 
expect more headroom on the discrete-memory T4 than on the
unified-memory M4, since unified memory shares one bandwidth
budget between CPU and GPU
(Section~\ref{rq2-hybrid-scheduling-conditions}).

\textbf{Design.} Native $r$ on the available hardware does
not span the regimes where the ratio can change the outcome, so
the GPU slowdown factor sweeps $r$ to find the hybrid applicability threshold (Section~\ref{methods-scheduling}); the queue-based
scheduler sees only task-completion timing, so the sleep is
sufficient at this granularity. Each workload/platform pair
first records $r_{\text{native}}$; if $r_{\text{native}}<1$,
the slowdown factor $s=1/r_{\text{native}}$ targets parity,
otherwise $s=1$. At the achieved ratio, \texttt{hybrid} is
paired against \texttt{vanilla} with the same slowdown
applied to the GPU side. The GPU path remains
\texttt{vanilla}, so the measured effect belongs to the
scheduler axis.

Native rows and slowed rows answer different questions. The
native-$r$ row runs with no added latency, so it is the only
row claiming a speedup realizable on this hardware. The slowed
rows distort $r$ to locate where adding CPU workers stops
helping, without claiming a gain on
this machine. Table~\ref{tab:hybrid-ratio-check} keeps native
$r$ and achieved $r$ in separate columns so the two are not
conflated.

\subsection{Calibration Cache Amortization
(RQ3)}\label{experiment-3-cache}

\textbf{Question.} When the unit of evaluation changes from
one grid to a related-checkpoint session, can the
unconditional calibration step amortize via 
cached-policy reuse?

\textbf{Motivation.} A LossLens-style
session~\cite{xie2025losslens} evaluates several related
checkpoints in succession. Hybrid scheduling's applicability
varies by machine and workload (RQ2), so a session cannot commit to hybrid statically. It must
calibrate, and calibration is allowed to select GPU-only as
a valid outcome. The question is whether the one-time
calibration cost is recovered over the session. RQ3 is a
scoped measurement: it uses one favorable regime and a
single session measurement per condition, so it supports or
refutes amortization for that case and does not establish a
general amortization claim across other workloads or machines.

\textbf{Hypothesis.} Amortization holds where Experiment~2
gives positive hybrid evidence and the per-checkpoint
hybrid saving exceeds the calibration cost amortized over
$N=4$.

\textbf{Design.} RQ3 filters Experiment~2 for a
(workload, regime) with hybrid affinity, meaning a supported
hybrid win. It prefers a native regime, where the win is a
practical-hardware claim, and falls back to a slowed regime,
where the win is a controlled demonstration. The selected
regime's slowdown is inherited as a single operating point and
applied uniformly to calibration, \texttt{vanilla}, and
\texttt{cached\_composed}, so the comparison is fair within
that regime. Each session contains $N=4$ checkpoints from one
training trajectory, drawn at evenly spaced epochs (the exact
spacing recorded with the asset manifest). The GPU-side
configuration is \texttt{rq3\_config}, selected by rule: the
RQ1 candidate with the highest supported speedup at the
workload/platform pair, taking the smaller $K$ on ties for
\texttt{vmap}-bearing candidates, and falling back to baseline
when no RQ1 candidate has a supported improvement. Two
conditions are compared: \texttt{vanilla} run once per
checkpoint, and \texttt{cached\_composed}, which calibrates on
checkpoint $0$ and reuses the selected tuple for all
checkpoints under \texttt{rq3\_config}. Amortization here means
\texttt{cached\_composed} has session speedup $>1.0$ after
calibration time is included; the verdict is labeled native or
slowed, and break-even $N^{*}$ follows
Section~\ref{methods-calibration}.

Two upstream outcomes shape RQ3. If RQ1 gives no supported
improvement at the selected pair, \texttt{rq3\_config} falls
back to baseline (Section~\ref{methods-calibration}), so RQ3
still runs but over the GPU-only evaluation path. If
Experiment~2 shows hybrid affinity in no regime, native or
slowed, there is no regime to inherit, and RQ3 reports
amortization refuted without running a session because
calibration cannot amortize anywhere. Within a selected regime,
calibration may still resolve to GPU-only, which also refutes
amortization. RQ3 therefore results in either a regime where
reuse amortizes or evidence that it does not.


\section{RESULTS}\label{results}

WIP

\section{DISCUSSION}\label{discussion}

WIP

\section{CONCLUSIONS}\label{conclusions}

WIP

\section{REFERENCES}\label{references}

\begin{thebibliography}{00}

\bibitem{goodfellow2015qualitative}
I. J. Goodfellow, O. Vinyals, and A. M. Saxe, Qualitatively
Characterizing Neural Network Optimization Problems.
International Conference on Learning Representations (ICLR),
2015. arXiv:1412.6544.
\url{https://arxiv.org/abs/1412.6544}

\bibitem{li2018losslandscape}
H. Li, Z. Xu, G. Taylor, C. Studer, and T. Goldstein,
Visualizing the Loss Landscape of Neural Nets. Advances in
Neural Information Processing Systems 31 (NeurIPS), 2018.
arXiv:1712.09913.
\url{https://papers.nips.cc/paper/7875-visualizing-the-loss-landscape-of-neural-nets}

\bibitem{xie2025losslens}
T. Xie, J. Chen, Y. Yang, C. Geniesse, G. Shi, A. J. Chaudhari, J. K.
Cava, M. W. Mahoney, T. Perciano, G. H. Weber, and R. Maciejewski,
LossLens: Diagnostics for Machine Learning Through Loss Landscape
Visual Analytics. IEEE Computer Graphics and Applications, 2025.
arXiv:2412.13321.
\url{https://doi.org/10.1109/MCG.2024.3509374}

\bibitem{losslens_repo}
T. Xie et al., LossLens: A Visual Analytics System for Diagnosing
Machine Learning Models via Loss Landscape Analysis. GitHub
repository, accessed May 22, 2026.
\url{https://github.com/VADERASU/loss-lens-CGA}

\bibitem{prakash2024adaptive}
A. Prakash, K. Wang, and R. Balestriero,
Visualizing Loss Landscapes with Adaptive Sampling and
Gradient Free Slicing. Brown University CS 237 Course
Proceedings, 2024.
\url{https://cs.brown.edu/courses/cs237/2025/proceedings2024.pdf}

\bibitem{polychronopoulos1987gss}
C. D. Polychronopoulos and D. J. Kuck, Guided Self-Scheduling: A
Practical Scheduling Scheme for Parallel Supercomputers. IEEE
Transactions on Computers, C-36(12), 1425--1439, 1987.
\url{https://doi.org/10.1109/TC.1987.5009495}

\bibitem{bora2022tinytasks}
S. Bora, B. Walker, and M. Fidler, The Tiny-Tasks Granularity
Trade-Off: Balancing overhead vs. performance in parallel systems.
arXiv:2202.11464, 2022.
\url{https://doi.org/10.48550/arXiv.2202.11464}

\bibitem{gallet2021heterogeneous}
B. Gallet and M. Gowanlock, Heterogeneous CPU-GPU Epsilon Grid
Joins: Static and Dynamic Work Partitioning Strategies. Data Science
and Engineering, 2021.
\url{https://doi.org/10.1007/s41019-020-00145-x}

\bibitem{khairy2020accelsim}
M. Khairy, Z. Shen, T. M. Aamodt, and T. Rogers, Accel-Sim: An
Extensible Simulation Framework for Validated GPU Modeling. Proceedings
of the 47th Annual International Symposium on Computer Architecture
(ISCA), 2020. \url{https://doi.org/10.1109/ISCA45697.2020.00047}

\bibitem{avalos2021pka}
C. Avalos, M. Khairy, R. Green, M. Payer, and T. Rogers,
Principal Kernel Analysis: A Tractable Methodology to Simulate
Scaled GPU Workloads. Proceedings of the 54th Annual IEEE/ACM
International Symposium on Microarchitecture (MICRO), 2021.
\url{https://doi.org/10.1145/3466752.3480100}

\bibitem{gabbay2022dnnmemory}
F. Gabbay, R. Lev Aharoni, and O. Schweitzer, Deep Neural Network
Memory Performance and Throughput Modeling and Simulation Framework.
Mathematics, 2022. \url{https://doi.org/10.3390/math10214144}

\bibitem{nvidia_cuda_guide}
NVIDIA, CUDA C++ Programming Guide, section 1.1, ``The Benefits
of Using GPUs''. NVIDIA documentation, accessed May 21, 2026.
\url{https://docs.nvidia.com/cuda/archive/11.1.1/cuda-c-programming-guide/index.html}

\bibitem{openvino_hetero}
Intel, Heterogeneous Execution (HETERO). OpenVINO documentation,
accessed May 22, 2026.
\url{https://docs.openvino.ai/2024/openvino-workflow/running-inference/inference-devices-and-modes/hetero-execution.html}

\bibitem{goldstein_loss_landscape_repo}
T. Goldstein et al. loss-landscape: Code for visualizing the loss
landscape of neural nets. GitHub repository, accessed May 21, 2026.
\url{https://github.com/tomgoldstein/loss-landscape}

\bibitem{pytorch_functional_call}
PyTorch, torch.func.functional\_call API documentation. PyTorch
documentation, accessed May 21, 2026.
\url{https://docs.pytorch.org/docs/stable/generated/torch.func.functional_call.html}

\bibitem{pytorch_vmap}
PyTorch, torch.vmap API documentation. PyTorch documentation,
accessed May 21, 2026.
\url{https://docs.pytorch.org/docs/stable/generated/torch.vmap.html}

\bibitem{pytorch_ensembling}
PyTorch, Model ensembling tutorial. PyTorch documentation,
accessed May 21, 2026.
\url{https://docs.pytorch.org/tutorials/intermediate/ensembling.html}

\bibitem{cho2014gru}
K. Cho, B. van Merrienboer, C. Gulcehre, D. Bahdanau, F. Bougares, H.
Schwenk, and Y. Bengio, Learning Phrase Representations using RNN
Encoder-Decoder for Statistical Machine Translation. Proceedings of
the 2014 Conference on Empirical Methods in Natural Language Processing
(EMNLP), 2014. arXiv:1406.1078.
\url{https://doi.org/10.3115/v1/D14-1179}

\bibitem{chung2014gru}
J. Chung, C. Gulcehre, K. Cho, and Y. Bengio, Empirical
Evaluation of Gated Recurrent Neural Networks on Sequence Modeling.
arXiv.org preprint arXiv:1412.3555, 2014.
\url{https://arxiv.org/abs/1412.3555}

\bibitem{holmes2019grnn}
C. Holmes, D. Mawhirter, Y. He, F. Yan, and B. Wu, GRNN:
Low-Latency and Scalable RNN Inference on GPUs. Proceedings of the
Fourteenth European Conference on Computer Systems (EuroSys), 2019.
\url{https://doi.org/10.1145/3302424.3303949}

\bibitem{georges2007benchmarking}
A. Georges, D. Buytaert, and L. Eeckhout, Statistically Rigorous
Java Performance Evaluation. Proceedings of the 22nd ACM SIGPLAN
Conference on Object-Oriented Programming Systems and Applications
(OOPSLA), 2007.
\url{https://doi.org/10.1145/1297027.1297033}

\bibitem{graham1982gprof}
S. L. Graham, P. B. Kessler, and M. K. McKusick, gprof: A Call
Graph Execution Profiler. SIGPLAN Symposium on Compiler Construction,
1982. \url{https://doi.org/10.1145/800230.806987}

\bibitem{ansel2014opentuner}
J. Ansel, S. Kamil, K. Veeramachaneni, J. Ragan-Kelley,
J. Bosboom, U.-M. O'Reilly, and S. Amarasinghe, OpenTuner: An
Extensible Framework for Program Autotuning. Proceedings of the
23rd International Conference on Parallel Architectures and
Compilation (PACT '14), pages 303--316, 2014.
\url{https://doi.org/10.1145/2628071.2628092}

\bibitem{whaley1998atlas}
R. C. Whaley and J. J. Dongarra, Automatically Tuned Linear
Algebra Software. Proceedings of the 1998 ACM/IEEE Conference
on Supercomputing (SC98), 1998. (Best Paper Award.)
\url{https://doi.org/10.1109/SC.1998.10004}

\end{thebibliography}

\section{APPENDICES (TODO)}\label{appendices}



\end{document}
